% RQ3: is the applied pacing delay objectively well calibrated?
%
% Merged from two candidate revisions.  What each contributed is recorded at the
% end of this block so neither line of reasoning is lost.
%
% Scope change from the baseline-comparison revision.  That table asked "is Ours
% better than another pacing policy" against Thermal LUT and CoDel-inspired
% arms.  Two problems retired it.  First, every collected cell for those arms
% stayed TBD.  Second, and decisively, a policy comparison is the wrong
% instrument for RQ3: RQ1(b) already establishes that pacing is necessary --
% the Admission-only and Pacing-only arms time out where the full controller
% does not -- so the remaining unanswered question is whether the delay the
% controller computes is the right *size*.  A baseline that paces differently
% cannot answer that; a requirement derived from the trace can.
%
% ---------------------------------------------------------------------------
% The measured requirement
% ---------------------------------------------------------------------------
% The deployed pacer (CaptureAvailablePacer.computePacingDelayMs) applies
%
%   d = ceil( max(0, Bhat + 2*Chat - max(0,T)) / 2 )
%
% with Bhat the backlog clock, Chat the reserved Draft duration, and T the
% window left on the newest committed capture's deadline.  The factor two is the
% two-Draft prospective horizon (the Draft that starts after the decision and
% the Draft of the capture the delay releases); the division by two is because
% one millisecond of delay both drains one millisecond of backlog and moves the
% next deadline out by one millisecond.  Verified against the workbook: the
% formula reproduces beforeAppliedDelayMs on 100% of recorded decisions in both
% conditions.
%
% The requirement keeps that condition and replaces only its two estimated
% inputs with their measured counterparts:
%
%   d* = ceil( max(0, B + 2*C - max(0,T)) / 2 )
%
%   B = RQ3Pacing.realBacklogMs        (max draftEndUptimeMs over unfinished
%                                       earlier Drafts, minus the decision time)
%   C = PacingReplay.draftSequenceDurationMs  (the Draft as it actually ran)
%   T = PacingReplay.beforeTimeToDeadlineMs
%
% This is deliberately not a different policy.  Because d and d* share the
% horizon model, the deadline, and the arithmetic, their difference isolates one
% thing only -- the delay attributable to the controller mispricing its own
% inputs -- rather than confounding estimator error with policy structure.
%
% ---------------------------------------------------------------------------
% Envelope, not tracking -- say this before a reviewer computes it
% ---------------------------------------------------------------------------
% d does not track d* per transition, and the manuscript must not imply that it
% does.  Over the transitions with a positive requirement:
%   Pearson r    -0.035 (12MP)  +0.134 (24MP)
%   Spearman     -0.048          +0.206
%   LS fit       d = -0.08 d* + 305 ms     d = 0.29 d* + 374 ms
%   d* quartiles  56 /  81 / 150 ms         78 / 187 / 342 ms
%   d  quartiles  26 / 323 / 474 ms          0 / 368 / 710 ms
% The requirement ranges over an order of magnitude while the applied delay is
% essentially uninformative about it -- in 12MP the relationship is now flat to
% slightly negative -- which is why the scatter reads as a vertical spread at
% each x rather than as points along the diagonal.
%
% That is arithmetic, not a defect.  Chat is
% CaptureAvailablePacingSession.getMaxDraftSequenceDurationMs -- the burst's
% observed *maximum* for the size bucket -- so it barely moves within a burst,
% while C is the actual duration of one Draft and moves a great deal.  A
% max-based reserve cannot correlate with a per-Draft actual.  The controller is
% therefore an envelope regulator, not a tracking one, which is the right shape
% for a deadline whose violation is user-visible and whose over-provisioning
% costs only bounded delay.
% It still responds at the coarse level, but only in 24MP: there d* <= 100 ms
% draws a median applied delay of 212 ms and d* > 300 ms draws 627 ms (n = 42
% and 43).  In 12MP the same buckets give 292 ms and 382 ms, and the upper one
% rests on four transitions, so the 12MP coarse response is not quotable.
%
% What is new in this data, and what the lower quartile of d says.  A quarter of
% the required transitions received 26 ms or less of delay in 12MP and none at
% all in 24MP: 18 of 79 and 43 of 140 required transitions were not paced.  The
% envelope is therefore not merely loose, it is absent on a substantial minority
% of the transitions that needed it, which is the reason coverage falls to 67.1%
% and 59.3%.  See the source note below -- this is concentrated in the second
% collection batch, and the manuscript must not describe the controller as
% holding an envelope without that qualification.
%
% Wording consequence.  "Calibration" is kept as the name of the *analysis*
% -- applied against required is exactly what it does -- but no caption, note,
% or sentence claims the controller computes a minimum, an optimum, or a match.
% The verbs are cover, clear, hold an envelope.  Claims to avoid: "the delay is
% calibrated to the requirement", "computes the minimum necessary delay",
% "tracks the requirement".
%
% Two variants of C:
%   C_mand = DynamicFunctionNode.durationMs + SecImageCodecNode ENCODING
%            durationMs + (draftSequenceDurationMs - workloadSequenceDurationMs)
%     The suffix admission can never skip, plus the non-node Draft overhead.
%     On the captures whose executed sequence demoted all the way to
%     DYNAMIC_FUNCTION>ENCODING (153 in 12MP, 276 in 24MP among the analyzed
%     transitions), DynamicFunction + Encoding equals workloadSequenceDurationMs
%     to the millisecond on every row.  d*_mand is a safety floor -- below it not
%     even mandatory work is protected.
%   C_exec = draftSequenceDurationMs, the configuration that actually ran.
%     d*_exec is the calibration target and is what "requirement" means below.
%
% The mandatory floor is kept out of the table as a column but it is NO LONGER a
% constant, and the change is the one finding here that runs against the
% controller.  It is positive on 29 of the 3,781 pooled transitions -- all in
% 24MP, none in 12MP -- and 14 of those 29 are breached: eleven applied no delay
% at all against a floor of 39--552 ms, and the remaining three fell short by a
% margin (4 vs 660, 627 vs 760, 764 vs 782 ms).  Under the previous run set the
% floor was positive on 6 transitions and
% breached on none, so a sentence of the form "not even the mandatory floor was
% ever violated" was available and now is not.  Do not write it.  The honest
% statement is that admission still carries most of the mandatory-work safety
% burden -- 3,752 of 3,781 transitions never put the floor above zero -- but
% that where the floor does bind, pacing covers it only about half the time.
%
% ---------------------------------------------------------------------------
% Why the surplus decomposition is exact
% ---------------------------------------------------------------------------
% Away from the max(0,.) clip -- transitions where both d > 0 and d* > 0 -- the
% two ceilings differ by at most half a millisecond and
%
%   d - d*  =  (Chat - C)  +  (Bhat - B)/2
%
% identically.  Those two terms are the Reserve and Backlog columns.  Checked on
% the 158 pooled transitions in that subset: residual mean 0.00 ms and
% |residual| < 1 ms on every row, i.e. rounding only.  Chat - C is
% PacingReplay.draftSequenceReserveErrorMs and Bhat - B is
% RQ3Pacing.backlogEstimateErrorMs, so both columns are also directly readable
% from the workbook without recomputation.
%
% This is the finding of the table.  The Reserve term carries the surplus in
% every row (+183 to +594 ms), while the Backlog term is an order of magnitude
% smaller and does not have a consistent sign: it is negative in seven of the
% eight rows (-30 to -108 ms), so the backlog clock is essentially unbiased and
% usually slightly optimistic, but 12MP Lv3-4 sits at +102 ms.  That row is the
% exception to state rather than to average away; it rests on 26 transitions and
% is the only place where the backlog clock, not the reserve, pays for a
% material part of the surplus.  The reserve behaviour is by construction:
% CaptureAvailablePacingSession.getMaxDraftSequenceDurationMs prices Chat at the
% burst's observed *maximum* Draft duration for the capture's size bucket,
% floored by the whole-Draft estimate.  A max-based reserve against a
% tail-latency deadline is the intended design; this table is what it costs, and
% it names the one estimator to attack if the cost is to come down.
%
% Cross-size contamination was investigated as an alternative explanation for
% the larger 24MP reserve term and rejected.  In the 24MP condition shots 1-2
% capture MP24 and shots 3+ MP12, so a heavy MP24 Draft can seed the reserve of
% a later MP12 capture (draftSequenceReserveCrossSizeContaminationMs > 0 on
% 18.3% of paced transitions).  But those transitions have a *smaller* median
% reserve term than the clean ones (347 vs 659 ms), so contamination is not the
% driver and is not claimed as one.
%
% ---------------------------------------------------------------------------
% Interpretation limit -- state it, do not hide it
% ---------------------------------------------------------------------------
% d* is computed on the realized trace, which pacing already shaped.  Removing
% the delay would deepen the backlog and lengthen later Drafts, so the true
% requirement of the unpaced trajectory is at or above d*.  The Surplus columns
% are therefore an UPPER bound on the delay that could have been removed, and
% the deficits are a LOWER bound.  Both directions are self-critical, which is
% what makes the estimator usable.
%
% The same limit is why "requirement" is never reported as "how often pacing was
% needed", and why the zero-requirement share is kept out of this table.  On the
% realized trace, 85.2% (12MP) and 79.4% (24MP) of paced transitions have
% d* = 0.  Read naively that is "most pacing was unnecessary", which is the
% closed-loop fallacy -- a regulator that holds the state away from the
% constraint makes every individual action look locally dispensable.  RQ1(b)
% supplies the counterfactual this local estimator cannot: with pacing removed,
% the same conditions time out.  Necessity is RQ1's question; RQ3 is scoped to
% size.  The two percentages appear in the caption of
% Figure~\ref{fig:rq3_pacing_calibration}, where the vertical-axis column they
% describe is visible and can be read together with its explanation.
%
% ---------------------------------------------------------------------------
% Data
% ---------------------------------------------------------------------------
% data/ablation_sampling/48U_metrics_12MP_normal_0803_{1,2}.xlsx and
% 48U_metrics_24MP_memory_0803_{1,2}.xlsx (ML implementation commit 99aae0a,
% exported by CaptureMetricsExcelExporter), both with admission and pacing
% enabled.  This is the same Full-arm source RQ1(a) now uses -- the balanced
% N = 10 per (condition, starting level) cell described in
% data/ablation_sampling/README.md -- so RQ1 and RQ3 finally describe one run
% set.  It replaces the single-workbook 0803 export the previous revision used,
% and the numbers move a long way; see "What changed with the run set" below.
%
% Run identity.  The two parts number their runs independently, so a run is keyed
% (part, runId) and no run id in these notes is comparable to one in the previous
% revision.  Eligibility is therefore stated as criteria, never as an id list.
%
% Run eligibility: complete 30-shot runs (RQ3Summary.isComplete30ShotRun),
% excluding runs that recorded a Capture Timeout.  In the 24MP condition a run
% captures MP24 on shots 1-2 and MP12 from shot 3, and from Lv5 the condition
% falls back to MP12 capture for the whole run, which is production behaviour and
% belongs to the condition; an MP12-from-shot-1 run below Lv5 is neither shape and
% is excluded.  This leaves 70 runs (ten per starting level) and 69 runs -- three
% short 24MP runs and one out-of-shape run are dropped, all four in Lv3 or Lv5.
% Transition eligibility: shots 2-30 with a recorded pacing decision and a
% complete prior-Draft timeline (pacingDecisionRecorded and
% realTraceCompleteBeforeDelay).  Two watchdog-truncated 24MP captures are dropped
% as transitions -- a truncated node has no clean executed-Draft duration -- while
% their runs are retained, consistent with RQ1(a) counting them among the
% zero-timeout runs.  Dropping those two runs whole instead leaves the
% required-transition set identical (140 positive, 59.3% covered).  This yields
% 1,920 and 1,861 analyzed transitions.
%
% Values quoted in the notes or held for prose (12MP / 24MP, pooled):
%   paced share            21.4 / 25.3 percent of transitions
%   requirement > 0        79 / 140 transitions
%   surplus P95            823 / 1,096 ms
%   burst surplus P50      2.91 / 1.40 s   (burst total delay 3.09 / 2.36 s)
%   required delay, d*>0   P50 81 / 187 ms   P95 301 / 663 ms
%   calibration error      P(< 0) 1.4 / 3.1 percent, P(= 0) 77.7 / 72.4
%   realized slack         P5 357 / 399 ms, min 8 / 14 ms, budget 7,000 ms
% Regenerate with scripts/rq3_calibration_metrics.py, which writes
% data/rq3/calibration/ and prints every number quoted in this file and in
% figures/fig_rq3_pacing_calibration.tex.
%
% ---------------------------------------------------------------------------
% What changed with the run set -- read before reusing any earlier sentence
% ---------------------------------------------------------------------------
% Against the previous single-workbook source (43 and 39 runs) coverage falls
% from 92.6% to 67.1% (12MP) and from 84.9% to 59.3% (24MP).  Two separable
% causes, and neither is a computation change: the estimator, the eligibility
% rules and the arithmetic all reproduce the previous table exactly when run on
% the previous source.
%   1. The second collection part.  Part 1 is close to the old export; part 2 is
%      a later batch in which the controller reserves much less (median Chat - C
%      of +186 and +135 ms against +263 and +365 ms in part 1) and its backlog
%      clock runs optimistic (median Bhat - B of -52 and -108 ms against +18 and
%      -27 ms).  Less over-reservation means less pacing -- 15.0% and 14.9% of
%      transitions paced in part 2 against 26.8% and 34.2% in part 1 -- and it
%      is where the unpaced-but-
%      required transitions live: 15 of 18 in 12MP and 40 of 43 in 24MP.  Read
%      per part, coverage is 81.1% / 54.8% (12MP) and 84.9% / 31.3% (24MP).
%      This asymmetry is large enough that it should be understood, and probably
%      explained in the paper, before these numbers are built on.
%   2. The RQ1(a) balancing trim, in 12MP only.  On the untrimmed run set the
%      12MP figures are 202 required transitions at 84.7% coverage; the trim to
%      ten runs per cell takes Lv3-4 from 117 required at 94.0% down to 31 at
%      77.4%.  data/ablation_sampling/README.md predeclares this: the trim scores
%      runs on reported outcome metrics and at Lv4 it drops the four runs that
%      retained all optional Draft work, which are exactly the runs whose heavy
%      Drafts generate large requirements.  The direction of that bias has to be
%      stated wherever the 12MP thermal-band rows are quoted.  In 24MP only two
%      runs were trimmed and the trim changes nothing (140 required, 59.3% both
%      before and after).
%
% Under-pacing does not track the tight realized margins: the smallest margin at
% an under-paced transition is 134 ms (12MP) and 36 ms (24MP), while the
% smallest margin overall is 8 ms and 14 ms.  The transitions that came closest
% to the deadline are therefore not the ones this estimator flags, which is
% consistent with d*_exec being a conservative target rather than a failure
% predictor.  Worth one sentence in prose.
%
% ---------------------------------------------------------------------------
% Layout, and what the merge kept from each revision
% ---------------------------------------------------------------------------
% Single-column float.  Ten candidate columns do not fit \columnwidth at a
% readable size, so the set was cut to what carries row-to-row variation:
% coverage with its denominator, the worst shortfall, the surplus, and the two
% terms it decomposes into.  Reserve and Backlog remain separate so their signs
% and relative magnitudes are explicit.  Waiting max is omitted because a
% one-off extreme is less stable and less direct than backlog P95.  Constants and
% single-number-per-condition quantities -- the mandatory floor, the per-burst
% surplus, realized slack -- moved to the caption and notes.
%
% Thermal bands.  Lv0-2 / Lv3-4 / Lv5-6, not the earlier two-band Lv0-3 /
% Lv4-6 split.  Both boundaries are product boundaries: the production guard of
% Section~\ref{sec:guard-limit} begins suppressing optional work at Lv4, and
% Lv5 is where the 24MP condition falls back to MP12 capture.  Three bands also
% resolve a pattern the two-band split hid.  Every column is monotone in thermal
% level, and identically so in both conditions:
%   coverage    42.4 -> 77.4 -> 100.0 %   and   30.8 -> 47.8 -> 89.1 %
%   surplus P50  191 ->  380 ->   407 ms  and    198 ->  255 ->  294 ms
%   backlog P95 71.0 -> 64.5 ->  56.6 %   and   75.4 -> 67.0 ->  52.3 %
% That is the policy characterisation this table exists to give: as thermal
% pressure rises the controller reserves more, and the extra reservation buys
% both higher coverage and a lighter regulated backlog.  On the balanced run set
% the trend is far stronger than it was -- the coverage span widens from
% 63.6-100.0 to 42.4-100.0 and from 72.7-88.6 to 30.8-89.1 -- because the run set
% now carries ten runs at every starting level instead of concentrating on
% Lv3-4.  The corollary is the one to state plainly: the controller holds its
% envelope under thermal pressure and loses it when the device is cool.  In the
% Lv0-2 band only 42.4% (12MP) and 30.8% (24MP) of the transitions that needed a
% delay received a sufficient one.
% Every band now rests on 15 or more required transitions (the thinnest is 12MP
% Lv5-6 at 15), so no cell is a three-point artefact; the denominators are still
% printed and should still travel with any quoted percentage.
%
% From the compact revision: the Lv0-4 / Lv5-6 resolution note, the run and
% session exclusions, and the discipline of stating that the target is
% trace-conditioned rather than a global counterfactual optimum.
% From the banded revision: the thermal-band rows, the surplus attribution, the
% worst-deficit and realized-slack anchors, and keeping the zero-requirement
% share out of the table.
% Not carried over: the P5/P50/P95 error triple, which mixes the under- and
% over-pacing tails of one distribution into a cell that has to be read in two
% directions; coverage plus worst deficit states the same two tails
% unambiguously.

\begin{table}[t]
  \centering
  \setlength{\abovecaptionskip}{2pt}%
  \setlength{\belowcaptionskip}{3pt}%
  \caption{RQ3 pacing calibration on the S26 Ultra}
  \label{tab:rq3_pacing_calibration}

  {\scriptsize
  \renewcommand{\arraystretch}{1.08}%
  \setlength{\arrayrulewidth}{0.35pt}%
  \setlength{\doublerulesep}{0.8pt}%
  \newsavebox{\rqthreecalbox}%
  \newcommand{\rqthreecaltabular}{%
  % Coverage carries "[83 of 140]" now that the 24MP denominators are three
  % digits, which no longer fits 30pt; the four millisecond columns hold at most
  % four glyphs, so the width comes out of Median and the table total is
  % unchanged.
  \begin{tabular}{|>{\centering\arraybackslash}m{25pt}%
                   |>{\centering\arraybackslash}m{23pt}%
                   ||>{\raggedleft\arraybackslash}m{36pt}%
                   |>{\raggedleft\arraybackslash}m{30pt}%
                   ||>{\raggedleft\arraybackslash}m{24pt}%
                   |>{\raggedleft\arraybackslash}m{30pt}%
                   |>{\raggedleft\arraybackslash}m{30pt}%
                   ||>{\raggedleft\arraybackslash}m{30pt}|}
    \hline\hline
    \multirow{2}{*}[-0.6ex]{\textbf{Cond.}} &
    \multirow{2}{*}[-0.6ex]{\makecell[c]{\textbf{Levels}\\\textbf{[\(n_{\mathrm{req}}\)]}}} &
    \multicolumn{2}{c||}{\textbf{Sufficiency}} &
    \multicolumn{3}{c||}{\textbf{Surplus decomposition (ms)}} &
    \multicolumn{1}{c|}{\makecell[c]{\textbf{Backlog}\\\textbf{state}}} \\
    \cline{3-8}
    & &
    \multicolumn{1}{c|}{\textbf{Coverage}} &
    \multicolumn{1}{c||}{\makecell[c]{\textbf{Worst}\\\textbf{deficit (ms)}}} &
    \multicolumn{1}{c|}{\textbf{Median}} &
    \multicolumn{1}{c|}{\makecell[c]{\textbf{Reserve}\\\textbf{term}}} &
    \multicolumn{1}{c||}{\makecell[c]{\textbf{Backlog}\\\textbf{term}}} &
    \multicolumn{1}{c|}{\makecell[c]{\textbf{P95}\\\textbf{(\% budget)}}} \\
    \hline\hline
    \multirow{4}{*}{\makecell[c]{12MP\\normal}}
      & \makecell[c]{Lv0--2\\{\scriptsize [33]}}
      & \makecell[r]{42.4\%\\{\scriptsize [14 of 33]}} & 205
      & 191 & $+309$ & $-108$ & 71.0\% \\
    \cline{2-8}
      & \makecell[c]{Lv3--4\\{\scriptsize [31]}}
      & \makecell[r]{77.4\%\\{\scriptsize [24 of 31]}} & 582
      & 380 & $+264$ & $+102$ & 64.5\% \\
    \cline{2-8}
      & \makecell[c]{Lv5--6\\{\scriptsize [15]}}
      & \makecell[r]{100.0\%\\{\scriptsize [15 of 15]}} & --
      & 407 & $+351$ & $0$ & 56.6\% \\
    \cline{2-8}
      & \makecell[c]{\textit{All}\\{\scriptsize [79]}}
      & \makecell[r]{67.1\%\\{\scriptsize [53 of 79]}} & 582
      & 357 & $+302$ & $-30$ & 65.4\% \\
    \hline
    \multirow{4}{*}{\makecell[c]{24MP\\memory\\pressure}}
      & \makecell[c]{Lv0--2\\{\scriptsize [39]}}
      & \makecell[r]{30.8\%\\{\scriptsize [12 of 39]}} & 552
      & 198 & $+183$ & $-34$ & 75.4\% \\
    \cline{2-8}
      & \makecell[c]{Lv3--4\\{\scriptsize [46]}}
      & \makecell[r]{47.8\%\\{\scriptsize [22 of 46]}} & 920
      & 255 & $+228$ & $-66$ & 67.0\% \\
    \cline{2-8}
      & \makecell[c]{Lv5--6\\{\scriptsize [55]}}
      & \makecell[r]{89.1\%\\{\scriptsize [49 of 55]}} & 195
      & 294 & $+594$ & $-60$ & 52.3\% \\
    \cline{2-8}
      & \makecell[c]{\textit{All}\\{\scriptsize [140]}}
      & \makecell[r]{59.3\%\\{\scriptsize [83 of 140]}} & 920
      & 257 & $+445$ & $-60$ & 69.5\% \\
    \hline\hline
  \end{tabular}}%
  \fittabcolsep{\rqthreecalbox}{\rqthreecaltabular}{\columnwidth}{16}%
  \rqthreecaltabular
  }
\end{table}
